% SEGMENT OLP-0385-S01
% Part: many-valued-logic
% Chapter: syntax-and-semantics
% Section: introduction

\documentclass[../../../include/open-logic-section]{subfiles}

\begin{document}

\olfileid{mvl}{syn}{int}

\olsection{அறிமுகம்}

செவ்வியல் தருக்கத்தில், $\lnot$, $\land$, $\lor$, $\lif$, $\liff$ என்ற
கூற்று இணைப்பிகளைப் பயன்படுத்தி !!{propositional variable}s இலிருந்து
கட்டப்படும் !!{formula}s ஐக் கையாள்கிறோம். செவ்வியல் தருக்கத்திற்குப்
பொருண்மையை வரையறுக்கும்போது, $\True$, $\False$ என்ற இரு மெய்மதிப்புகளைப்
பயன்படுத்துகிறோம். !!{propositional variable}s ஐ
!!a{valuation}~$\pAssign{v}$ இல் பொருள்கோள் கொள்கிறோம்; அது அந்த
!!{propositional variable}s க்கு $\True$, $\False$ என்ற மெய்மதிப்புகளை
ஒதுக்குகிறது. பின்னர் எந்த !!{valuation} உம் எந்த
!!{formula}~$!A$ க்கும் ஒரு மெய்மதிப்பு~$\pValue{v}(!A)$ ஐத்
தீர்மானிக்கிறது; மேலும் !!a{valuation}~$\pAssign{v}$ இல் !!a{formula}
நிறைவுறுத்தப்படுகிறது, அதாவது $\pSat{v}{!A}$ என்பது, அப்போதும் அப்போது
மட்டுமே, $\pValue{v}(!A) = \True$.

பல்மதிப்புத் தருக்கங்கள், $\True$, $\False$ என்ற இரண்டை மட்டும் விடக்
கூடுதலான மெய்மதிப்புகளை அனுமதிப்பதன் மூலம் செவ்வியல் இருமதிப்புத்
தருக்கத்தைப் பொதுமைப்படுத்துகின்றன. எனவே பல்மதிப்புத் தருக்கத்தில்,
!!a{valuation}~$\pAssign{v}$ என்பது ஒவ்வொரு !!{propositional
variable}~$p$ க்கும் கிடைக்கக்கூடிய மெய்மதிப்புகளின் ஒரு வீச்சிலிருந்து
ஒரு மதிப்பை ஒதுக்கும் சார்பு. அனுமதிக்கப்பட்ட மெய்மதிப்புகளின் கணத்தைப்
பொதுவாக~$V$ என்போம். $V = \{\True, \False\}$ ஆகவும், இணைப்பிகளுக்கான
பழக்கமான வரையறைமுறை மெய்மை அட்டவணைகளைக் கொண்டு
மெய்மதிப்பு~$\pValue{v}(!A)$ கணிக்கப்படுவதாகவும் உள்ள பல்மதிப்புத்
தருக்கமே செவ்வியல் தருக்கம்.

கூடுதல் மெய்மதிப்புகளைச் சேர்த்தவுடன், ``மற்றும்,'' ``அல்லது,'' ``அன்று,''
``எனில்---அப்போது'' என்று வாசிக்கும் இணைப்பிகளுக்கு~$\pValue{v}(!A)$ ஐ
எவ்வாறு கணிப்பது என்பதற்குப் பல இயல்பான வாய்ப்புகள் கிடைக்கின்றன. எனவே ஒரு
பல்மதிப்புத் தருக்கத்தை மெய்மதிப்புகளின் கணம் மட்டும் தீர்மானிப்பதில்லை;
ஒவ்வொரு இணைப்பிக்கும் நாம் தேர்ந்தெடுக்கும் \emph{வாய்மைச் சார்புகளும்}
தீர்மானிக்கின்றன. ஆனால் பல்மதிப்புத் தருக்கம்~$\Log L$ ஒன்றுக்கு இவை
தேர்ந்தெடுக்கப்பட்டவுடன், செவ்வியல் தருக்கத்தில் போலவே மதிப்பீடு
மெய்மதிப்பு~$\pValue{v}(!A)[\Log L]$ ஐத் தனித்தவாறு தீர்மானிக்கிறது.
செவ்வியல் தருக்கத்தைப் போலவே பல்மதிப்புத் தருக்கங்களும்
\emph{வாய்மைச் சார்புகளால் தீர்மானிக்கப்படுகின்றன}.

% SOURCE CORRECTION TA-MVL-001: repair agreement in “these semantic building blocks” and read “role,” not “rule,” of True.
இந்தப் பொருண்மை அடிக்கட்டுகளைக் கொண்டு, மெய்மம், பின்விளைவு,
நிறைவுறுத்தத்தக்க தன்மை ஆகிய பொருண்மைக் கருத்துகளுக்கான ஒப்புமைகளை நாம்
வரையறுக்கலாம். செவ்வியல் தருக்கத்தில், ஒவ்வொரு~$\pAssign{v}$ க்கும் அதன்
மெய்மதிப்பு $\pValue{v}(!A) = \True$ எனில், !!a{formula} ஒரு மெய்மம்.
பல்மதிப்புத் தருக்கத்தில் இதையும் சிறிது பொதுமைப்படுத்த வேண்டும். முதலில்,
மெய்மதிப்புகளின் கணம்~$V$ ஆனது~$\True$ ஐக் கொண்டிருக்க வேண்டும் என்ற
கட்டாயமில்லை. எடுத்துக்காட்டாக, சில பல்மதிப்புத் தருக்கங்கள் $0$ க்கும்~$1$
க்கும் இடையிலுள்ள எல்லா விகிதமுறு எண்கள் போன்ற எண்களைத் தமது
மெய்மதிப்புகளின் கணமாகப் பயன்படுத்துகின்றன. அப்போது $1$ வழக்கமாக~$\True$
இன் இடத்தை வகிக்கிறது. வேறு தருக்கங்களில், ஒரே ஒரு மெய்மதிப்பு மட்டுமன்றி
பல மெய்மதிப்புகள் அவ்விடத்தை வகிக்கின்றன. ஆகவே ஒவ்வொரு பல்மதிப்புத்
தருக்கத்திற்கும் \emph{நியமிக்கப்பட்ட மதிப்புகள்} அடங்கிய ஒரு கணம்~$V^+$
இருக்க வேண்டும் என்று கோருகிறோம். பின்னர் !!a{valuation}~$\pAssign{v}$
இல் !!a{formula} நிறைவுறுத்தப்படுகிறது, அதாவது
$\pSat{v}{!A}[\Log L]$ என்பது, அப்போதும் அப்போது மட்டுமே,
$\pValue{v}(!A)[\Log L] \in V^+$ எனக் கூறலாம். ஒவ்வொரு~$\pAssign{v}$
க்கும் $\pValue{v}(!A) \in V^+$ எனில், !!a{formula}~$!A$ அந்தத்
தருக்கத்தின் மெய்மம், அதாவது $\Entails[\Log L] !A$. இறுதியாக,
~$\Gamma$ இல் உள்ள எல்லா !!{formula}s ஐயும் நிறைவுறுத்தும் ஒவ்வொரு
!!{valuation} உம்~$!A$ ஐயும் நிறைவுறுத்தினால், !!{formula}s இன்
அக்கணம்~$!A$ ஐப் பின்விளைவாகக் கொண்டுள்ளது, அதாவது
$\Gamma \Entails[\Log L] !A$ என்கிறோம்.

\end{document}
